% =============================================================
% 00_intro.tex  導入：試験の全体像と本の使い方
% =============================================================
\chapter{ITパスポートの全体像と、本書の歩き方}

\section{ITパスポートとは何か}

ITパスポート試験は、情報処理推進機構（IPA）が実施する国家試験のうち、もっとも入り口に位置づけられた試験です。
エンジニアだけでなく、あらゆる社会人が「IT とビジネスを共通言語で話すための最低限」を身につけることを目的にしています。

\definitionbx{この試験の位置づけ}{
IT を「使う側」「発注する側」「一緒に働く側」の人が、共通のことばで意思疎通できるようになるための試験。
エンジニア向けの専門試験ではなく、\termtag{全社会人向けの IT 基礎教養}に近い。
}

\pitfall{「専門試験ではないのに、なぜこんなに用語が出てくるのか」}{
ITパスポートは「広く浅く」の広さがかなり広いのが特徴です。
経営・法律・プロジェクト管理・ネットワーク・セキュリティ・AI まで、守備範囲が一気に広がります。
\textbf{これは網羅の罠です}。暗記で全部覚えようとすると必ず息切れします。
本書では、\termtag{用語同士のつながり}を先に押さえ、個別の単語をそのつながりに貼り付けていく読み方を提案します。
}

\section{試験の全体像}

試験は大きく 3 分野に分かれています。それぞれが独立しているように見えて、実は\textbf{関係の順番}があります。

\figurebox{ITパスポート 3 分野の関係図}{fig:three-fields}{%
  % 図: ITパスポート3分野の関係図
% 用途: % 図: ITパスポート3分野の関係図
% 用途: % 図: ITパスポート3分野の関係図
% 用途: \input{assets/diagrams/three_fields}
\begin{tikzpicture}[node distance=8mm and 8mm]

  \node[keynode, minimum width=28mm] (strategy)
    {\textbf{ストラテジ系}\\\footnotesize 経営・法務\\\footnotesize 目的を決める};
  \node[keynode, minimum width=28mm, right=of strategy] (mgmt)
    {\textbf{マネジメント系}\\\footnotesize 開発・運用\\\footnotesize どう回すか};
  \node[keynode, minimum width=28mm, right=of mgmt] (tech)
    {\textbf{テクノロジ系}\\\footnotesize HW/NW/DB/SEC\\\footnotesize 何でできているか};

  \draw[arr2] (strategy) -- node[above, font=\scriptsize]{要件} (mgmt);
  \draw[arr2] (mgmt) -- node[above, font=\scriptsize]{設計} (tech);

  \node[subnode, below=6mm of mgmt, minimum width=92mm, text width=88mm]
    (read) {\textbf{左が目的／右が手段}。ストラテジ→マネジメント→テクノロジの順に抽象度が下がる。};

  \begin{scope}[on background layer]
    \node[groupbg, fit=(strategy)(mgmt)(tech), inner sep=10pt] {};
  \end{scope}

\end{tikzpicture}

\begin{tikzpicture}[node distance=8mm and 8mm]

  \node[keynode, minimum width=28mm] (strategy)
    {\textbf{ストラテジ系}\\\footnotesize 経営・法務\\\footnotesize 目的を決める};
  \node[keynode, minimum width=28mm, right=of strategy] (mgmt)
    {\textbf{マネジメント系}\\\footnotesize 開発・運用\\\footnotesize どう回すか};
  \node[keynode, minimum width=28mm, right=of mgmt] (tech)
    {\textbf{テクノロジ系}\\\footnotesize HW/NW/DB/SEC\\\footnotesize 何でできているか};

  \draw[arr2] (strategy) -- node[above, font=\scriptsize]{要件} (mgmt);
  \draw[arr2] (mgmt) -- node[above, font=\scriptsize]{設計} (tech);

  \node[subnode, below=6mm of mgmt, minimum width=92mm, text width=88mm]
    (read) {\textbf{左が目的／右が手段}。ストラテジ→マネジメント→テクノロジの順に抽象度が下がる。};

  \begin{scope}[on background layer]
    \node[groupbg, fit=(strategy)(mgmt)(tech), inner sep=10pt] {};
  \end{scope}

\end{tikzpicture}

\begin{tikzpicture}[node distance=8mm and 8mm]

  \node[keynode, minimum width=28mm] (strategy)
    {\textbf{ストラテジ系}\\\footnotesize 経営・法務\\\footnotesize 目的を決める};
  \node[keynode, minimum width=28mm, right=of strategy] (mgmt)
    {\textbf{マネジメント系}\\\footnotesize 開発・運用\\\footnotesize どう回すか};
  \node[keynode, minimum width=28mm, right=of mgmt] (tech)
    {\textbf{テクノロジ系}\\\footnotesize HW/NW/DB/SEC\\\footnotesize 何でできているか};

  \draw[arr2] (strategy) -- node[above, font=\scriptsize]{要件} (mgmt);
  \draw[arr2] (mgmt) -- node[above, font=\scriptsize]{設計} (tech);

  \node[subnode, below=6mm of mgmt, minimum width=92mm, text width=88mm]
    (read) {\textbf{左が目的／右が手段}。ストラテジ→マネジメント→テクノロジの順に抽象度が下がる。};

  \begin{scope}[on background layer]
    \node[groupbg, fit=(strategy)(mgmt)(tech), inner sep=10pt] {};
  \end{scope}

\end{tikzpicture}
%
}

\comparisonbx{3 分野の読み方}{%
\comptable{分野 & 一言で言うと & 典型的な問い}{
\textbf{ストラテジ系} & 何のために作るか & 経営戦略・法務・財務 \\
\textbf{マネジメント系} & どう作り、どう回すか & 開発工程・プロジェクト管理・サービス運用 \\
\textbf{テクノロジ系} & 何でできているか & ハード・ネットワーク・DB・セキュリティ
}
}

\advice{%
用語を覚える前に、まず\textbf{どの分野の話か}を決めます。
「法務」と聞いたらストラテジ、「納期」と聞いたらマネジメント、「暗号」と聞いたらテクノロジ。
この 3 秒の仕分けを反射でできるようにすることが、合格への最短ルートです。
}

\section{合格のための勉強法（本書の使い方）}

\subsection{やり切るための 3 つの原則}

\begin{enumerate}
  \item \textbf{一冊を往復する}：一度で覚えきろうとしない。本書を 3 周する前提で、1 周目はスピード重視、2 周目で比較、3 周目で確認問題を解く。
  \item \textbf{関係図を先に}：各章冒頭の図を先に見て、後から本文で肉付けする。単語単独で覚えない。
  \item \textbf{ひっかけの理由まで読む}：解説の\textbf{「なぜ他の選択肢が間違いか」}まで言葉にできれば本番で迷わない。
\end{enumerate}

\subsection{章の読み方}

各節は、読者がつまずきやすい場所を起点に、次の流れで書かれています。

\begin{itemize}
  \item まず何に困るか（\termtag{つまずきポイント}）
  \item 直感的なたとえ
  \item 正しい定義
  \item 間違えやすい比較
  \item 試験でどう問われるか
  \item 1 問確認
\end{itemize}

\onelinesummary{%
ITパスポートは「IT を使う側」の共通言語試験。
3 分野は\textbf{「目的→手段→部品」}の順で読む。
本書は、用語を点で覚えるのではなく、\termtag{関係図で覚える}ことを目的とする。
}

\checkquestion{ITパスポート試験について、最も適切な記述はどれか。}
  {IT エンジニアを目指す学生のための、専門家育成試験である。}
  {IT 技術者以外を含む、すべての社会人を想定した IT 基礎教養試験である。}
  {ネットワーク設計の実務試験である。}
  {プログラミング言語の記述力を問う試験である。}
  {B}
  {ITパスポートは「IT を使う側の共通言語」を問う試験です。A は基本情報技術者以降に近く、C はネットワークスペシャリスト、D は言語別試験の発想で、いずれもスコープが狭すぎます。}

\supportqr{about}{本書の使い方（サポートサイト）}
